\section{Experimental Setup}
\label{sec:experiments}

We design experiments to answer the following research questions:

\textbf{RQ1 (H-E1):} Can temporal dynamic attention be implemented and trained stably?

\textbf{RQ2 (H-M1):} Does reducing temporal steps ($T$) decrease FLOPs while maintaining perplexity?

\textbf{RQ3 (H-M2):} Does attention entropy decrease across temporal steps (convergence)?

\textbf{RQ4 (H-C1):} Does FLOP reduction scale with sequence length?

\subsection{Dataset}

We evaluate on \textbf{WikiText-103}, a standard language modeling benchmark with $\sim$103M training tokens.

\subsection{Model Architecture}

\begin{table}[h]
\centering
\begin{tabular}{@{}ll@{}}
\toprule
\textbf{Component} & \textbf{Configuration} \\
\midrule
Model & Decoder-only Transformer \\
Layers & 6 \\
Hidden dimension & 512 \\
Attention heads & 8 \\
Temporal steps $T$ & \{1, 2, 3\} (varied) \\
\bottomrule
\end{tabular}
\caption{Model architecture configuration.}
\label{tab:architecture}
\end{table}

\textbf{Variants:} baseline ($T$=1), temporal\_t3 ($T$=3), temporal\_t2 ($T$=2), cached\_kv\_t3 (ablation).

\subsection{Implementation Details}

\textbf{Framework:} PyTorch 2.0 \\
\textbf{Training:} AdamW, lr=$1 \times 10^{-4}$, batch size 32, 10 epochs \\
\textbf{FLOP measurement:} \texttt{thop} library

\subsection{Evaluation Metrics}

\begin{itemize}
    \item \textbf{Perplexity (PPL):} $\exp(\text{cross-entropy loss})$
    \item \textbf{FLOPs (G):} Floating-point operations in billions
    \item \textbf{Attention Entropy:} $H = -\sum p \log p$
    \item \textbf{FLOP Reduction:} $(\text{baseline} - \text{method}) / \text{baseline} \times 100\%$
\end{itemize}
